\documentclass[9pt,a4paper]{article}
\usepackage[utf8]{inputenc}
\usepackage[spanish]{babel}
\usepackage{amsmath,amssymb,amsthm}
\usepackage{multicol}
\usepackage{geometry}
\usepackage{enumitem}
\usepackage{fancyhdr}
\usepackage{xcolor}
\usepackage{tcolorbox}
\usepackage{booktabs}
\usepackage{array}
\usepackage{graphicx}

% Márgenes reducidos
\geometry{left=0.7cm,right=0.7cm,top=0.9cm,bottom=0.9cm}

% Configuración de listas compactas
\setlist{nosep,leftmargin=*,topsep=0pt,partopsep=0pt}

% Comandos personalizados
\newcommand{\E}{\mathbb{E}}
\newcommand{\Var}{\text{Var}}
\newcommand{\Cov}{\text{Cov}}
\newcommand{\sd}{\tilde{s}}
\newcommand{\barra}{\overline}
\newcommand{\R}{\texttt{R: }}

% Colores
\definecolor{seccolor}{RGB}{0,70,120}
\definecolor{boxcolor}{RGB}{235,245,255}
\definecolor{alertcolor}{RGB}{255,245,235}
\definecolor{excolor}{RGB}{245,255,245}

% Secciones compactas
\usepackage{titlesec}
\titleformat{\section}{\normalfont\bfseries\color{seccolor}}{\thesection.}{0.3em}{}
\titlespacing*{\section}{0pt}{5pt}{2pt}
\titleformat{\subsection}{\normalfont\bfseries\footnotesize\color{seccolor}}{\thesubsection.}{0.3em}{}
\titlespacing*{\subsection}{0pt}{3pt}{1pt}
\titleformat{\subsubsection}{\normalfont\bfseries\scriptsize\color{seccolor!80}}{\thesubsubsection.}{0.3em}{}
\titlespacing*{\subsubsection}{0pt}{2pt}{1pt}

% Cajas
\newtcolorbox{formulabox}[1][]{colback=boxcolor,colframe=seccolor,boxrule=0.4pt,left=2pt,right=2pt,top=1pt,bottom=1pt,arc=1pt,#1}
\newtcolorbox{alertbox}[1][]{colback=alertcolor,colframe=orange!70!black,boxrule=0.4pt,left=2pt,right=2pt,top=1pt,bottom=1pt,arc=1pt,#1}
\newtcolorbox{exbox}[1][]{colback=excolor,colframe=green!50!black,boxrule=0.4pt,left=2pt,right=2pt,top=1pt,bottom=1pt,arc=1pt,#1}

\pagestyle{fancy}
\fancyhf{}
\fancyhead[L]{\footnotesize\textbf{Formulario de Estadística - Resolución de Problemas}}
\fancyhead[R]{\footnotesize Pág. \thepage/6}
\renewcommand{\headrulewidth}{0.3pt}
\setlength{\columnsep}{0.4cm}

\begin{document}
\begin{multicols}{2}

%===============================================
\section{PROBABILIDAD}
%===============================================

\subsection{Definiciones básicas}
\textbf{Espacio muestral} $\Omega$: conjunto de todos los resultados posibles.\\
\textbf{Suceso}: $A\subseteq\Omega$. \quad \textbf{Suceso seguro}: $\Omega$. \quad \textbf{Suceso imposible}: $\emptyset$.\\
\textbf{Complementario}: $A^c=\Omega\setminus A$

\textbf{Axiomas de Kolmogorov}:
\begin{enumerate}[leftmargin=*]
\item $P(A)\geq 0$ para todo $A$
\item $P(\Omega)=1$
\item Si $A\cap B=\emptyset$: $P(A\cup B)=P(A)+P(B)$
\end{enumerate}

\begin{formulabox}
\textbf{Propiedades}:\\
$P(\emptyset)=0$ \quad $P(A^c)=1-P(A)$\\
$P(A\cup B)=P(A)+P(B)-P(A\cap B)$\\
$P(A\cup B\cup C)=P(A)+P(B)+P(C)-P(A\cap B)-P(A\cap C)-P(B\cap C)+P(A\cap B\cap C)$
\end{formulabox}

\subsection{Probabilidad Condicionada}
\begin{formulabox}
$P(A|B)=\frac{P(A\cap B)}{P(B)}$ \quad (si $P(B)>0$)
\end{formulabox}

\textbf{Regla del producto}: $P(A\cap B)=P(A|B)\cdot P(B)=P(B|A)\cdot P(A)$

\textbf{Regla de la cadena}:\\
$P(A_1\cap\cdots\cap A_n)=P(A_1)\cdot P(A_2|A_1)\cdot P(A_3|A_1\cap A_2)\cdots$

\subsection{Independencia de sucesos}
$A$ y $B$ son \textbf{independientes} si y solo si:
\begin{itemize}[leftmargin=*]
\item $P(A\cap B)=P(A)\cdot P(B)$
\item $P(A|B)=P(A)$ (cuando $P(B)>0$)
\item $P(B|A)=P(B)$ (cuando $P(A)>0$)
\end{itemize}

\begin{alertbox}
\textbf{¡IMPORTANTE!}\\
$A,B$ disjuntos ($A\cap B=\emptyset$) $\neq$ $A,B$ independientes\\
Si $A,B$ disjuntos con $P(A),P(B)>0$, entonces son \textbf{dependientes}.
\end{alertbox}

\textbf{Independencia mutua}: $A_1,\ldots,A_n$ indep. si para todo subconjunto $I\subseteq\{1,\ldots,n\}$:\\
$P(\bigcap_{i\in I}A_i)=\prod_{i\in I}P(A_i)$

\subsection{Teorema de la Probabilidad Total}
Si $\{A_1,\ldots,A_n\}$ es \textbf{partición} de $\Omega$ (disjuntos 2 a 2 y $\cup A_i=\Omega$):
\begin{formulabox}
$P(B)=\sum_{i=1}^{n}P(B|A_i)\cdot P(A_i)$
\end{formulabox}

\subsection{Teorema de Bayes}
\begin{formulabox}
$P(A_k|B)=\frac{P(B|A_k)\cdot P(A_k)}{\sum_{i=1}^{n}P(B|A_i)\cdot P(A_i)}=\frac{P(B|A_k)\cdot P(A_k)}{P(B)}$
\end{formulabox}

\begin{exbox}
\textbf{Ej Bayes}: Test médico: sensib. 99\%, especif. 95\%, preval. 0.5\%.\\
$P(\text{enf}|+)=\frac{0.99\cdot 0.005}{0.99\cdot 0.005+0.05\cdot 0.995}\approx 9\%$
\end{exbox}

\subsection{Combinatoria}
\begin{center}
\begin{tabular}{@{}lll@{}}
\textbf{Tipo} & \textbf{Fórmula} & \textbf{Descripción}\\
\hline
Permutaciones & $P_n=n!$ & Ordenar $n$ objetos\\
Variaciones & $V_n^k=\frac{n!}{(n-k)!}$ & Elegir $k$ de $n$ con orden\\
Var. con rep. & $VR_n^k=n^k$ & Elegir $k$ de $n$ con rep. y orden\\
Combinaciones & $\binom{n}{k}=\frac{n!}{k!(n-k)!}$ & Elegir $k$ de $n$ sin orden\\
Comb. con rep. & $CR_n^k=\binom{n+k-1}{k}$ & $k$ de $n$ con rep. sin orden\\
Perm. con rep. & $\frac{n!}{n_1!\cdots n_k!}$ & $n$ objetos con repeticiones
\end{tabular}
\end{center}

\textbf{Propiedades binomiales}:\\
$\binom{n}{0}=\binom{n}{n}=1$, \quad $\binom{n}{k}=\binom{n}{n-k}$\\
$\binom{n}{k}=\binom{n-1}{k-1}+\binom{n-1}{k}$ (Pascal)

%===============================================
\section{VARIABLES ALEATORIAS}
%===============================================

\subsection{Definición}
\textbf{Variable aleatoria (v.a.)}: $X:\Omega\to\mathbb{R}$, asigna valor numérico a cada resultado.\\
\textbf{Dominio}: $D_X=\{x\in\mathbb{R}: P(X=x)>0\}$ (discretas)

\subsection{V.A. Discreta}
$D_X$ finito o numerable.

\textbf{Función de probabilidad}: $P_X(x)=P(X=x)$\\
Propiedades: $P_X(x)\geq 0$, \quad $\sum_{x\in D_X}P_X(x)=1$

\textbf{Función de distribución}: 
\begin{formulabox}
$F_X(x)=P(X\leq x)=\sum_{x_i\leq x}P_X(x_i)$
\end{formulabox}
Propiedades: $0\leq F_X(x)\leq 1$, no decreciente, $\lim_{x\to-\infty}F_X(x)=0$, $\lim_{x\to+\infty}F_X(x)=1$

$P(a<X\leq b)=F_X(b)-F_X(a)$

\subsection{V.A. Continua}
$D_X$ es un intervalo o unión de intervalos.

\textbf{Función de densidad}: $f_X(x)$\\
Propiedades: $f_X(x)\geq 0$, \quad $\int_{-\infty}^{\infty}f_X(x)dx=1$\\
\textbf{Nota}: $P(X=x)=0$ para todo $x$ (probabilidad puntual nula)

\textbf{Función de distribución}:
\begin{formulabox}
$F_X(x)=P(X\leq x)=\int_{-\infty}^{x}f_X(t)dt$ \quad $f_X(x)=F_X'(x)$
\end{formulabox}

$P(a<X<b)=P(a\leq X\leq b)=\int_a^b f_X(x)dx=F_X(b)-F_X(a)$

\subsection{Esperanza (valor esperado)}
\begin{formulabox}
\textbf{Discreta}: $\E(X)=\mu_X=\sum_{x\in D_X} x\cdot P_X(x)$\\[3pt]
\textbf{Continua}: $\E(X)=\mu_X=\int_{-\infty}^{\infty}x\cdot f_X(x)dx$
\end{formulabox}

\textbf{Esperanza de función}:\\
$\E(g(X))=\sum_x g(x)P_X(x)$ \quad o \quad $\E(g(X))=\int g(x)f_X(x)dx$

\textbf{Propiedades}:
\begin{itemize}[leftmargin=*]
\item $\E(c)=c$ (constante)
\item $\E(aX+b)=a\E(X)+b$ (linealidad)
\item $\E(X+Y)=\E(X)+\E(Y)$ (siempre, aunque dependientes)
\item Si $X,Y$ indep.: $\E(XY)=\E(X)\E(Y)$
\end{itemize}

\subsection{Varianza y desviación típica}
\begin{formulabox}
$\Var(X)=\sigma_X^2=\E[(X-\mu_X)^2]=\E(X^2)-[\E(X)]^2$\\[3pt]
$\sigma_X=+\sqrt{\Var(X)}$ \quad (desviación típica)
\end{formulabox}

\textbf{Propiedades}:
\begin{itemize}[leftmargin=*]
\item $\Var(X)\geq 0$, \quad $\Var(c)=0$
\item $\Var(aX+b)=a^2\Var(X)$
\item Si $X,Y$ indep.: $\Var(X+Y)=\Var(X)+\Var(Y)$
\item General: $\Var(X+Y)=\Var(X)+\Var(Y)+2\Cov(X,Y)$
\end{itemize}

\subsection{Covarianza y correlación}
\begin{formulabox}
$\Cov(X,Y)=\E[(X-\mu_X)(Y-\mu_Y)]=\E(XY)-\E(X)\E(Y)$\\[3pt]
$\rho_{XY}=\frac{\Cov(X,Y)}{\sigma_X\sigma_Y}$ \quad (coef. correlación, $-1\leq\rho\leq 1$)
\end{formulabox}

Si $X,Y$ indep. $\Rightarrow$ $\Cov(X,Y)=0$ (pero no al revés)

\subsection{Cuantiles}
El \textbf{cuantil $q$} de $X$ es $x_q$ tal que $F_X(x_q)=P(X\leq x_q)=q$

\textbf{Percentiles}: cuantil expresado en porcentaje (ej: percentil 95 = cuantil 0.95)

\textbf{Mediana}: cuantil 0.5

%===============================================
\section{DISTRIBUCIONES DISCRETAS}
%===============================================

\subsection{Bernoulli $Ber(p)$}
Experimento con dos resultados: éxito (1) con prob. $p$, fracaso (0) con prob. $q=1-p$.
\begin{formulabox}
$D_X=\{0,1\}$, \quad $P(X=x)=p^x(1-p)^{1-x}$\\
$\E(X)=p$, \quad $\Var(X)=p(1-p)=pq$
\end{formulabox}

\subsection{Binomial $B(n,p)$}
$X=$ número de éxitos en $n$ ensayos Bernoulli independientes con $P(\text{éxito})=p$.
\begin{formulabox}
$D_X=\{0,1,\ldots,n\}$\\
$P(X=k)=\binom{n}{k}p^k(1-p)^{n-k}$\\
$\E(X)=np$, \quad $\Var(X)=np(1-p)$
\end{formulabox}

\textbf{Propiedad}: $B(n,p)=\sum_{i=1}^n Ber_i(p)$ (suma de Bernoulli indep.)

\R\texttt{dbinom(k,size=n,prob=p)} (prob.), \texttt{pbinom(k,n,p)} (distrib.)

\begin{exbox}
\textbf{Ej}: Lanzar moneda ($p=0.5$) 10 veces. $P(X\geq 8)$?\\
$P(X\geq 8)=P(X=8)+P(X=9)+P(X=10)=\binom{10}{8}0.5^{10}+\ldots$\\
\R \texttt{1-pbinom(7,10,0.5)} o \texttt{pbinom(7,10,0.5,lower.tail=FALSE)}
\end{exbox}

\subsection{Geométrica $Ge(p)$}
$X=$ número de \textbf{fracasos} antes del primer éxito.
\begin{formulabox}
$D_X=\{0,1,2,\ldots\}$\\
$P(X=k)=(1-p)^k\cdot p=q^k p$\\
$F(k)=P(X\leq k)=1-q^{k+1}$\\
$\E(X)=\frac{1-p}{p}=\frac{q}{p}$, \quad $\Var(X)=\frac{1-p}{p^2}=\frac{q}{p^2}$
\end{formulabox}

\begin{alertbox}
\textbf{Falta de memoria}: $P(X>k+j|X\geq j)=P(X>k)$\\
Tras $j$ fracasos, la prob. de $k$ fracasos más es la misma que al inicio.
\end{alertbox}

\R\texttt{dgeom(k,prob=p)}, \texttt{pgeom(k,p)} (nota: cuenta fracasos)

\subsection{Binomial Negativa $BN(r,p)$}
$X=$ número de fracasos antes del $r$-ésimo éxito.
\begin{formulabox}
$D_X=\{0,1,2,\ldots\}$\\
$P(X=k)=\binom{k+r-1}{r-1}(1-p)^k p^r$\\
$\E(X)=\frac{r(1-p)}{p}$, \quad $\Var(X)=\frac{r(1-p)}{p^2}$
\end{formulabox}

\textbf{Nota}: $BN(1,p)=Ge(p)$

\R\texttt{dnbinom(k,size=r,prob=p)}, \texttt{pnbinom(k,r,p)}

\subsection{Poisson $Po(\lambda)$}
$X=$ número de eventos en intervalo con tasa media $\lambda$.
\begin{formulabox}
$D_X=\{0,1,2,\ldots\}$\\
$P(X=k)=\frac{\lambda^k e^{-\lambda}}{k!}$\\
$\E(X)=\lambda$, \quad $\Var(X)=\lambda$
\end{formulabox}

\textbf{Aproximación de Binomial}: Si $n\geq 20$, $p\leq 0.05$, $np\leq 10$:\\
$B(n,p)\approx Po(np)$

\textbf{Proceso de Poisson}: Si tasa es $\lambda$ por unidad de tiempo, entonces en intervalo de longitud $t$: $X_t\sim Po(\lambda t)$

\R\texttt{dpois(k,lambda)}, \texttt{ppois(k,lambda)}

\begin{exbox}
\textbf{Ej}: Llegan 3 clientes/hora en promedio. $P(\text{más de 5 en 2h})$?\\
$X_{2h}\sim Po(6)$. $P(X>5)=1-P(X\leq 5)$\\
\R \texttt{1-ppois(5,6)} $\approx 0.446$
\end{exbox}

\subsection{Hipergeométrica $H(m,n,k)$}
Urna con $m$ bolas blancas y $n$ rojas. Se extraen $k$ sin reposición.\\
$X=$ número de bolas blancas extraídas.
\begin{formulabox}
$D_X=\{\max(0,k-n),\ldots,\min(m,k)\}$\\
$P(X=x)=\frac{\binom{m}{x}\binom{n}{k-x}}{\binom{m+n}{k}}$\\
$\E(X)=\frac{km}{m+n}$\\
$\Var(X)=k\cdot\frac{m}{m+n}\cdot\frac{n}{m+n}\cdot\frac{m+n-k}{m+n-1}$
\end{formulabox}

\R\texttt{dhyper(x,m,n,k)}, \texttt{phyper(x,m,n,k)}

%===============================================
\section{DISTRIBUCIONES CONTINUAS}
%===============================================

\subsection{Uniforme $U(a,b)$}
Todos los valores en $(a,b)$ igualmente probables.
\begin{formulabox}
$f(x)=\begin{cases}\frac{1}{b-a} & \text{si }a<x<b\\0 & \text{otro caso}\end{cases}$\\[5pt]
$F(x)=\begin{cases}0 & x\leq a\\\frac{x-a}{b-a} & a<x<b\\1 & x\geq b\end{cases}$\\[5pt]
$\E(X)=\frac{a+b}{2}$, \quad $\Var(X)=\frac{(b-a)^2}{12}$
\end{formulabox}

\R\texttt{dunif(x,min=a,max=b)}, \texttt{punif}, \texttt{qunif}, \texttt{runif}

\subsection{Exponencial $Exp(\lambda)$}
Tiempo entre eventos de proceso Poisson con tasa $\lambda$.
\begin{formulabox}
$f(x)=\begin{cases}\lambda e^{-\lambda x} & x\geq 0\\0 & x<0\end{cases}$\\[5pt]
$F(x)=\begin{cases}1-e^{-\lambda x} & x\geq 0\\0 & x<0\end{cases}$\\[5pt]
$\E(X)=\frac{1}{\lambda}$, \quad $\Var(X)=\frac{1}{\lambda^2}$
\end{formulabox}

\begin{alertbox}
\textbf{Falta de memoria}: $P(X>s+t|X>s)=P(X>t)=e^{-\lambda t}$
\end{alertbox}

\R\texttt{dexp(x,rate=lambda)}, \texttt{pexp}, \texttt{qexp}

\begin{exbox}
\textbf{Ej}: Bombilla con $E(X)=10000$h. $P(X>12000|X>1000)$?\\
$\lambda=1/10000$. Por falta de memoria: $P(X>11000)=e^{-11000/10000}\approx 0.333$
\end{exbox}

\subsection{Normal (Gaussiana) $N(\mu,\sigma)$}
\begin{formulabox}
$f(x)=\frac{1}{\sigma\sqrt{2\pi}}e^{-\frac{(x-\mu)^2}{2\sigma^2}}$, \quad $x\in\mathbb{R}$\\[3pt]
$\E(X)=\mu$, \quad $\Var(X)=\sigma^2$
\end{formulabox}

\textbf{Normal estándar} $Z\sim N(0,1)$: $\phi(z)=f_Z(z)$, $\Phi(z)=F_Z(z)$

\textbf{Tipificación}: Si $X\sim N(\mu,\sigma)$, entonces $Z=\frac{X-\mu}{\sigma}\sim N(0,1)$

\textbf{Propiedades de simetría}:
\begin{itemize}[leftmargin=*]
\item $\Phi(-z)=1-\Phi(z)$
\item $P(Z>z)=1-\Phi(z)$
\item $P(-z<Z<z)=2\Phi(z)-1$
\item $z_q=-z_{1-q}$ (cuantiles simétricos)
\end{itemize}

\textbf{Propiedades de combinación}:
\begin{itemize}[leftmargin=*]
\item $aX+b\sim N(a\mu+b, |a|\sigma)$
\item Si $X_i\sim N(\mu_i,\sigma_i)$ indep.: $\sum X_i\sim N(\sum\mu_i,\sqrt{\sum\sigma_i^2})$
\end{itemize}

\textbf{Cuantiles importantes de $N(0,1)$}:
\begin{center}
\begin{tabular}{c|ccccc}
$q$ & 0.90 & 0.95 & 0.975 & 0.99 & 0.995\\
\hline
$z_q$ & 1.282 & 1.645 & 1.960 & 2.326 & 2.576
\end{tabular}
\end{center}

\R\texttt{dnorm(x,mean=mu,sd=sigma)}, \texttt{pnorm}, \texttt{qnorm}

\begin{exbox}
\textbf{Ej}: $X\sim N(100,15)$. $P(X>120)$?\\
$P(X>120)=P(Z>\frac{120-100}{15})=P(Z>1.33)=1-\Phi(1.33)$\\
\R \texttt{1-pnorm(120,100,15)} $\approx 0.0912$
\end{exbox}

\subsection{Chi-cuadrado $\chi^2_n$}
Si $Z_1,\ldots,Z_n\stackrel{iid}{\sim}N(0,1)$: $\sum_{i=1}^n Z_i^2\sim\chi^2_n$
\begin{formulabox}
$\E(\chi^2_n)=n$, \quad $\Var(\chi^2_n)=2n$
\end{formulabox}

\textbf{Resultado clave}: Si $X_1,\ldots,X_n\stackrel{iid}{\sim}N(\mu,\sigma)$:\\
$\frac{(n-1)\tilde{S}^2_X}{\sigma^2}\sim\chi^2_{n-1}$

\textbf{Nota}: $\chi^2$ NO es simétrica. Para IC de varianza, los cuantiles van ``al revés''.

\R\texttt{dchisq(x,df=n)}, \texttt{pchisq}, \texttt{qchisq}

\subsection{t de Student $t_n$}
Si $Z\sim N(0,1)$ y $V\sim\chi^2_n$ independientes: $T=\frac{Z}{\sqrt{V/n}}\sim t_n$
\begin{formulabox}
$\E(t_n)=0$ (si $n>1$), \quad $\Var(t_n)=\frac{n}{n-2}$ (si $n>2$)
\end{formulabox}

\textbf{Resultado clave}: Si $X_1,\ldots,X_n\stackrel{iid}{\sim}N(\mu,\sigma)$:\\
$T=\frac{\bar{X}-\mu}{\tilde{S}_X/\sqrt{n}}\sim t_{n-1}$

\textbf{Propiedades}:
\begin{itemize}[leftmargin=*]
\item Simétrica respecto a 0
\item $t_{n,q}=-t_{n,1-q}$
\item Para $n\geq 30$: $t_n\approx N(0,1)$
\item Más "pesada en colas" que $N(0,1)$
\end{itemize}

\R\texttt{dt(x,df=n)}, \texttt{pt}, \texttt{qt}

\subsection{F de Fisher $F_{m,n}$}
Si $U\sim\chi^2_m$, $V\sim\chi^2_n$ indep.: $F=\frac{U/m}{V/n}\sim F_{m,n}$

\textbf{Resultado clave}: Si $X_i\stackrel{iid}{\sim}N(\mu_1,\sigma_1)$, $Y_j\stackrel{iid}{\sim}N(\mu_2,\sigma_2)$ indep.:\\
$\frac{\tilde{S}^2_X/\sigma_1^2}{\tilde{S}^2_Y/\sigma_2^2}\sim F_{n_1-1,n_2-1}$

\textbf{Propiedad}: $1/F_{m,n}\sim F_{n,m}$

\R\texttt{df(x,df1,df2)}, \texttt{pf}, \texttt{qf}

%===============================================
\section{MUESTREO Y ESTIMACIÓN}
%===============================================

\subsection{Conceptos básicos}
\textbf{Población}: conjunto de todos los elementos de interés.\\
\textbf{Muestra}: subconjunto de la población.\\
\textbf{M.a.s.}: $X_1,\ldots,X_n$ v.a. i.i.d. con misma distribución que la población.\\
\textbf{Parámetro}: característica numérica de la población ($\mu$, $\sigma$, $p$).\\
\textbf{Estadístico}: función de la muestra (es una v.a.).\\
\textbf{Estimador}: estadístico usado para estimar un parámetro.

\subsection{Tipos de muestreo}
\begin{itemize}[leftmargin=*]
\item \textbf{Con reposición}: cada elemento puede ser seleccionado varias veces
\item \textbf{Sin reposición}: cada elemento se selecciona una sola vez
\item \textbf{Sistemático}: seleccionar cada $k$-ésimo elemento de una lista
\item \textbf{Estratificado}: dividir población en estratos y muestrear cada uno
\item \textbf{Por conglomerados}: seleccionar grupos completos al azar
\end{itemize}

\begin{alertbox}
\textbf{Regla}: Si $N\geq 1000n$, muestreo con/sin reposición son equivalentes.
\end{alertbox}

\textbf{Factor de corrección para población finita} (muestreo sin reposición, $n$ grande respecto a $N$):
\begin{formulabox}
$\sigma_{\bar{X}}=\frac{\sigma}{\sqrt{n}}\cdot\sqrt{\frac{N-n}{N-1}}$\quad
$\sigma_{\hat{p}}=\sqrt{\frac{p(1-p)}{n}}\cdot\sqrt{\frac{N-n}{N-1}}$
\end{formulabox}

\subsection{Estadísticos muestrales}
\begin{formulabox}
\textbf{Media muestral}: $\bar{X}=\frac{1}{n}\sum_{i=1}^n X_i$\\[5pt]
\textbf{Varianza muestral}: $\tilde{S}^2_X=\frac{1}{n-1}\sum_{i=1}^n(X_i-\bar{X})^2$\\[5pt]
\textbf{Cuasivarianza}: $S^2_X=\frac{1}{n}\sum_{i=1}^n(X_i-\bar{X})^2=\frac{n-1}{n}\tilde{S}^2_X$\\[5pt]
\textbf{Proporción muestral}: $\hat{p}=\frac{\text{núm. éxitos}}{n}=\frac{\sum X_i}{n}$
\end{formulabox}

\textbf{Fórmula alternativa para varianza}:\\
$\tilde{S}^2_X=\frac{1}{n-1}\left(\sum X_i^2-n\bar{X}^2\right)=\frac{n}{n-1}\left(\frac{\sum X_i^2}{n}-\bar{X}^2\right)$

\subsection{Propiedades de la media muestral}
Si $X_1,\ldots,X_n$ es m.a.s. de $X$ con $\E(X)=\mu$, $\Var(X)=\sigma^2$:
\begin{formulabox}
$\E(\bar{X})=\mu$ \quad (insesgado)\\
$\Var(\bar{X})=\frac{\sigma^2}{n}$\\
$\sigma_{\bar{X}}=\frac{\sigma}{\sqrt{n}}$ \quad (error estándar)
\end{formulabox}

Si además $X\sim N(\mu,\sigma)$: $\bar{X}\sim N\left(\mu,\frac{\sigma}{\sqrt{n}}\right)$

\subsection{Teorema Central del Límite (TCL)}
\begin{formulabox}
Si $X_1,\ldots,X_n$ m.a.s. con $\E(X_i)=\mu$, $\Var(X_i)=\sigma^2$, entonces para $n$ grande:\\[3pt]
$\bar{X}\approx N\left(\mu,\frac{\sigma}{\sqrt{n}}\right)$ \quad equiv. \quad $\frac{\bar{X}-\mu}{\sigma/\sqrt{n}}\approx N(0,1)$
\end{formulabox}

\textbf{Regla práctica}: $n\geq 30$ suele ser suficiente.

\subsection{Propiedades de estimadores}
\begin{itemize}[leftmargin=*]
\item \textbf{Insesgado}: $\E(\hat{\theta})=\theta$ (sesgo $B(\hat{\theta})=\E(\hat{\theta})-\theta=0$)
\item \textbf{Consistente}: $\hat{\theta}\xrightarrow{P}\theta$ cuando $n\to\infty$
\item \textbf{Eficiente}: mínima varianza entre insesgados
\item \textbf{Error estándar}: $\sigma_{\hat{\theta}}=\sqrt{\Var(\hat{\theta})}$
\item \textbf{ECM}: $\text{ECM}(\hat{\theta})=\Var(\hat{\theta})+[B(\hat{\theta})]^2$
\end{itemize}

\begin{alertbox}
\textbf{¡OJO con la varianza!}\\
$\tilde{S}^2_X=\frac{\sum(X_i-\bar{X})^2}{n-1}$ es \textbf{insesgado}: $\E(\tilde{S}^2_X)=\sigma^2$\\
$S^2_X=\frac{\sum(X_i-\bar{X})^2}{n}$ tiene \textbf{sesgo}: $\E(S^2_X)=\frac{n-1}{n}\sigma^2$ (sesgo $=-\frac{\sigma^2}{n}\to 0$)\\
Para $n$ pequeño ($\leq 30$), usar $\tilde{S}^2_X$. Para $n$ grande, ambas valen.
\end{alertbox}

\textbf{Estimadores insesgados eficientes} (para población normal):
\begin{itemize}[leftmargin=*]
\item $\bar{X}$ para $\mu$
\item $\tilde{S}^2_X$ para $\sigma^2$
\item $\hat{p}$ para $p$ (Bernoulli)
\end{itemize}

\subsection{Distribuciones muestrales importantes}
Si $X_1,\ldots,X_n\stackrel{iid}{\sim}N(\mu,\sigma)$:
\begin{formulabox}
$\bar{X}\sim N\left(\mu,\frac{\sigma}{\sqrt{n}}\right)$\\[3pt]
$\frac{\bar{X}-\mu}{\sigma/\sqrt{n}}\sim N(0,1)$\\[3pt]
$\frac{\bar{X}-\mu}{\tilde{S}_X/\sqrt{n}}\sim t_{n-1}$\\[3pt]
$\frac{(n-1)\tilde{S}^2_X}{\sigma^2}\sim\chi^2_{n-1}$
\end{formulabox}

%===============================================
\section{INTERVALOS DE CONFIANZA}
%===============================================

\subsection{Definición}
Un \textbf{IC del $(1-\alpha)\cdot 100\%$} para $\theta$ es $(A,B)$ tal que $P(A<\theta<B)=1-\alpha$
\begin{itemize}[leftmargin=*]
\item $1-\alpha$: nivel de confianza (típicamente 0.95)
\item $\alpha$: nivel de significación
\item $\alpha/2$: probabilidad en cada cola (IC bilateral)
\end{itemize}

\subsection{IC para $\mu$ de población normal, $\sigma$ conocida}
\begin{formulabox}
$\left(\bar{X}-z_{1-\alpha/2}\frac{\sigma}{\sqrt{n}},\; \bar{X}+z_{1-\alpha/2}\frac{\sigma}{\sqrt{n}}\right)$\\[3pt]
Abreviado: $\bar{X}\pm z_{1-\alpha/2}\frac{\sigma}{\sqrt{n}}$
\end{formulabox}

\textbf{Amplitud}: $A=2z_{1-\alpha/2}\frac{\sigma}{\sqrt{n}}$

\textbf{Tamaño muestral para amplitud máxima $A_0$}:\\
$n\geq\left(\frac{2z_{1-\alpha/2}\sigma}{A_0}\right)^2$

\subsection{IC para $\mu$ de población normal, $\sigma$ desconocida}
\begin{formulabox}
$\bar{X}\pm t_{n-1,1-\alpha/2}\frac{\tilde{S}_X}{\sqrt{n}}$
\end{formulabox}

\R\texttt{t.test(X,conf.level=0.95)\$conf.int}

\subsection{IC para $\mu$ de población cualquiera ($n$ grande)}
Si $n\geq 40$, por TCL:
\begin{formulabox}
$\bar{X}\pm z_{1-\alpha/2}\frac{\tilde{S}_X}{\sqrt{n}}$
\end{formulabox}

\subsection{IC para proporción $p$}

\textbf{Método de Laplace} (aproximación normal):\\
Condiciones: $n\geq 100$, $n\hat{p}\geq 10$, $n(1-\hat{p})\geq 10$
\begin{formulabox}
$\hat{p}\pm z_{1-\alpha/2}\sqrt{\frac{\hat{p}(1-\hat{p})}{n}}$
\end{formulabox}
\R\texttt{binom.approx(x,n,conf.level)} (paquete epitools)

\textbf{Método de Wilson} (mejor para muestras pequeñas):
\begin{formulabox}
$\frac{\hat{p}+\frac{z^2}{2n}\pm z\sqrt{\frac{\hat{p}(1-\hat{p})}{n}+\frac{z^2}{4n^2}}}{1+\frac{z^2}{n}}$
\end{formulabox}
donde $z=z_{1-\alpha/2}$\\
\R\texttt{binom.wilson(x,n,conf.level)} (paquete epitools)

\textbf{Método de Clopper-Pearson} (exacto):\\
\R\texttt{binom.exact(x,n,conf.level)} (paquete epitools)

\textbf{Tamaño muestral para proporción} (peor caso $\hat{p}=0.5$):
\begin{formulabox}
$n\geq\left\lceil\frac{z_{1-\alpha/2}^2}{A^2}\right\rceil$ donde $A$ es semi-amplitud deseada
\end{formulabox}

\begin{exbox}
\textbf{Ej}: IC 95\% con error máximo 0.05: $n\geq\lceil\frac{1.96^2}{0.1^2}\rceil=385$
\end{exbox}

\subsection{IC para varianza $\sigma^2$ (población normal)}
\begin{formulabox}
$\left(\frac{(n-1)\tilde{S}^2_X}{\chi^2_{n-1,1-\alpha/2}},\;\frac{(n-1)\tilde{S}^2_X}{\chi^2_{n-1,\alpha/2}}\right)$
\end{formulabox}

\textbf{IC para $\sigma$}: tomar raíces cuadradas de los extremos.

\R\texttt{varTest(X,conf.level)\$conf.int} (paquete EnvStats)

\subsection{IC para diferencia de medias (muestras indep.)}
\textbf{Varianzas conocidas}:
$(\bar{X}_1-\bar{X}_2)\pm z_{1-\alpha/2}\sqrt{\frac{\sigma_1^2}{n_1}+\frac{\sigma_2^2}{n_2}}$

\textbf{Varianzas desconocidas iguales}:
$(\bar{X}_1-\bar{X}_2)\pm t_{n_1+n_2-2,1-\alpha/2}\cdot S_p\sqrt{\frac{1}{n_1}+\frac{1}{n_2}}$\\
donde $S_p^2=\frac{(n_1-1)\tilde{S}_1^2+(n_2-1)\tilde{S}_2^2}{n_1+n_2-2}$

\textbf{Varianzas desconocidas distintas (Welch)}:
$(\bar{X}_1-\bar{X}_2)\pm t_{\nu,1-\alpha/2}\sqrt{\frac{\tilde{S}_1^2}{n_1}+\frac{\tilde{S}_2^2}{n_2}}$

%===============================================
\section{CONTRASTE DE HIPÓTESIS}
%===============================================

\subsection{Conceptos básicos}
\begin{itemize}[leftmargin=*]
\item $H_0$: \textbf{Hipótesis nula} (lo que se contrasta, status quo)
\item $H_1$: \textbf{Hipótesis alternativa} (lo que queremos demostrar)
\item \textbf{Estadístico de contraste}: función de la muestra para tomar decisión
\item \textbf{p-valor}: prob. de obtener resultado tan extremo o más, si $H_0$ es cierta
\item \textbf{Región crítica}: valores del estadístico que llevan a rechazar $H_0$
\end{itemize}

\textbf{Método de los 6 pasos}:
\begin{enumerate}[leftmargin=*]
\item Establecer $H_0$ y $H_1$
\item Fijar nivel de significación $\alpha$
\item Seleccionar estadístico de contraste
\item Calcular valor del estadístico con los datos
\item Calcular p-valor
\item Tomar decisión: rechazar $H_0$ si p-valor $<\alpha$
\end{enumerate}

\textbf{Tipos de hipótesis alternativa}:
\begin{itemize}[leftmargin=*]
\item \textbf{Unilateral derecha}: $H_1:\theta>\theta_0$
\item \textbf{Unilateral izquierda}: $H_1:\theta<\theta_0$
\item \textbf{Bilateral}: $H_1:\theta\neq\theta_0$
\end{itemize}

\begin{alertbox}
\textbf{Regla de decisión con p-valor}:\\
$p<0.05$: Rechazar $H_0$ (significativo *)\\
$p<0.01$: Rechazar $H_0$ (muy significativo **)\\
$p<0.001$: Rechazar $H_0$ (altamente significativo ***)\\
$p>0.1$: No rechazar $H_0$\\
$0.05\leq p\leq 0.1$: Zona de penumbra (más datos)\\[3pt]
\textbf{¡OJO!} p-valor $\neq$ $P(H_0\text{ cierta})$
\end{alertbox}

\subsection{Tipos de errores}
\begin{center}
\begin{tabular}{@{}l|cc@{}}
& $H_0$ cierta & $H_0$ falsa\\
\hline
Aceptar $H_0$ & Correcto & Error Tipo II ($\beta$)\\
Rechazar $H_0$ & Error Tipo I ($\alpha$) & Correcto (Potencia $1-\beta$)
\end{tabular}
\end{center}
$\alpha$: nivel de significación\\
$1-\beta$: potencia del test

\subsection{Contraste para $\mu$ ($\sigma$ conocida) - Z-test}
$H_0:\mu=\mu_0$\\
Estadístico: $Z=\frac{\bar{X}-\mu_0}{\sigma/\sqrt{n}}\sim N(0,1)$ bajo $H_0$

\begin{center}
\begin{tabular}{@{}l|l@{}}
$H_1$ & p-valor\\
\hline
$\mu>\mu_0$ & $P(Z\geq z_0)=1-\Phi(z_0)$\\
$\mu<\mu_0$ & $P(Z\leq z_0)=\Phi(z_0)$\\
$\mu\neq\mu_0$ & $2P(Z\geq|z_0|)=2(1-\Phi(|z_0|))$
\end{tabular}
\end{center}

\subsection{Contraste para $\mu$ ($\sigma$ desconocida) - t-test}
$H_0:\mu=\mu_0$\\
Estadístico: $T=\frac{\bar{X}-\mu_0}{\tilde{S}_X/\sqrt{n}}\sim t_{n-1}$ bajo $H_0$

\begin{center}
\begin{tabular}{@{}l|l@{}}
$H_1$ & p-valor\\
\hline
$\mu>\mu_0$ & $P(t_{n-1}\geq t_0)$\\
$\mu<\mu_0$ & $P(t_{n-1}\leq t_0)$\\
$\mu\neq\mu_0$ & $2P(t_{n-1}\geq|t_0|)$
\end{tabular}
\end{center}

\R\texttt{t.test(X,mu=mu0,alternative="two.sided/greater/less")}

\subsection{Contraste para proporción $p$}
$H_0:p=p_0$

\textbf{Exacto (binomial)}:\\
Estadístico: $X\sim B(n,p_0)$ bajo $H_0$\\
\R\texttt{binom.test(x,n,p=p0,alternative=...)}

\textbf{Aproximado ($n$ grande)}:\\
Estadístico: $Z=\frac{\hat{p}-p_0}{\sqrt{p_0(1-p_0)/n}}\sim N(0,1)$\\
\R\texttt{prop.test(x,n,p=p0,alternative=...)}

\subsection{Contraste para $\sigma^2$ - $\chi^2$-test}
$H_0:\sigma^2=\sigma_0^2$\\
Estadístico: $\chi^2=\frac{(n-1)\tilde{S}^2_X}{\sigma_0^2}\sim\chi^2_{n-1}$ bajo $H_0$

\begin{center}
\begin{tabular}{@{}l|l@{}}
$H_1$ & p-valor\\
\hline
$\sigma>\sigma_0$ & $P(\chi^2_{n-1}\geq\chi^2_0)$\\
$\sigma<\sigma_0$ & $P(\chi^2_{n-1}\leq\chi^2_0)$\\
$\sigma\neq\sigma_0$ & $2\min\{P(\leq\chi^2_0),P(\geq\chi^2_0)\}$
\end{tabular}
\end{center}

\R\texttt{sigma.test(X,sigma=sigma0,...)} (paquete TeachingDemos)

\subsection{Contraste para dos medias independientes}
$H_0:\mu_1=\mu_2$

\textbf{Varianzas iguales}:
\begin{formulabox}
$T=\frac{\bar{X}_1-\bar{X}_2}{S_p\sqrt{\frac{1}{n_1}+\frac{1}{n_2}}}\sim t_{n_1+n_2-2}$\\[5pt]
$S_p^2=\frac{(n_1-1)\tilde{S}_1^2+(n_2-1)\tilde{S}_2^2}{n_1+n_2-2}$
\end{formulabox}

\textbf{Varianzas distintas (Welch)}:
\begin{formulabox}
$T=\frac{\bar{X}_1-\bar{X}_2}{\sqrt{\frac{\tilde{S}_1^2}{n_1}+\frac{\tilde{S}_2^2}{n_2}}}\sim t_\nu$
\end{formulabox}
(grados de libertad $\nu$ aproximados por fórmula de Welch)

\R\texttt{t.test(X,Y,var.equal=TRUE/FALSE,alternative=...)}

\subsection{Contraste para dos proporciones}
$H_0:p_1=p_2$
\begin{formulabox}
$Z=\frac{\hat{p}_1-\hat{p}_2}{\sqrt{\hat{p}(1-\hat{p})(\frac{1}{n_1}+\frac{1}{n_2})}}\sim N(0,1)$\\[5pt]
donde $\hat{p}=\frac{x_1+x_2}{n_1+n_2}$ (proporción combinada)
\end{formulabox}

\R\texttt{prop.test(c(x1,x2),c(n1,n2),alternative=...)}

\textbf{Test exacto de Fisher}: \texttt{fisher.test(tabla,...)}

\subsection{Contraste para dos varianzas - F-test}
$H_0:\sigma_1^2=\sigma_2^2$ (equiv. $\frac{\sigma_1^2}{\sigma_2^2}=1$)\\
Estadístico: $F=\frac{\tilde{S}_1^2}{\tilde{S}_2^2}\sim F_{n_1-1,n_2-1}$ bajo $H_0$

\begin{center}
\begin{tabular}{@{}l|l@{}}
$H_1$ & p-valor\\
\hline
$\sigma_1>\sigma_2$ & $P(F_{n_1-1,n_2-1}\geq f_0)$\\
$\sigma_1<\sigma_2$ & $P(F_{n_1-1,n_2-1}\leq f_0)$\\
$\sigma_1\neq\sigma_2$ & $2\min\{P(\leq f_0),P(\geq f_0)\}$
\end{tabular}
\end{center}

\R\texttt{var.test(X,Y,alternative=...)}

\textbf{Test no paramétrico (si normalidad dudosa)}:\\
\R\texttt{fligner.test(list(X,Y))} (solo bilateral)

\subsection{Muestras emparejadas}
Para datos pareados $(X_i,Y_i)$: calcular diferencias $D_i=X_i-Y_i$.

$H_0:\mu_1=\mu_2$ equiv. $H_0:\mu_D=0$
\begin{formulabox}
$T=\frac{\bar{D}}{\tilde{S}_D/\sqrt{n}}\sim t_{n-1}$ bajo $H_0$
\end{formulabox}

\textbf{IC para $\mu_D=\mu_1-\mu_2$}: $\bar{D}\pm t_{n-1,1-\alpha/2}\frac{\tilde{S}_D}{\sqrt{n}}$

\R\texttt{t.test(X,Y,paired=TRUE,alternative=...)}

\begin{alertbox}
\textbf{¿Cuándo usar emparejadas?}\\
Mismos individuos antes/después, gemelos, medidas repetidas.\\
El diseño experimental se fija \textbf{ANTES} de recoger datos.
\end{alertbox}

%===============================================
\section{BONDAD DE AJUSTE}
%===============================================

\subsection{Test $\chi^2$ de Pearson}
$H_0$: La muestra proviene de la distribución teórica especificada.

\begin{formulabox}
$\chi^2=\sum_{i=1}^{k}\frac{(O_i-E_i)^2}{E_i}\sim\chi^2_{k-1-r}$
\end{formulabox}
donde:
\begin{itemize}[leftmargin=*]
\item $O_i$: frecuencia observada de la clase $i$
\item $E_i=n\cdot p_i$: frecuencia esperada
\item $k$: número de clases
\item $r$: número de parámetros estimados de los datos
\end{itemize}

\textbf{Condiciones de aplicación}:
\begin{itemize}[leftmargin=*]
\item $n\geq 30$
\item $E_i\geq 5$ para todo $i$ (si no, agrupar clases)
\end{itemize}

p-valor $=P(\chi^2_{g.l.}\geq\chi^2_0)$ donde g.l.$=k-1-r$

\R\texttt{chisq.test(x,p=probabilidades)}\\
Para simular: \texttt{chisq.test(...,simulate.p.value=TRUE,B=5000)}

\begin{exbox}
\textbf{Ej}: Dado lanzado 120 veces. Observaciones: 20,22,17,18,19,24.\\
$H_0$: dado equilibrado ($p_i=1/6$). $E_i=120/6=20$.\\
$\chi^2=\frac{(20-20)^2}{20}+\frac{(22-20)^2}{20}+\cdots=1.7$\\
p-valor $=P(\chi^2_5\geq 1.7)\approx 0.89$. No rechazar $H_0$.
\end{exbox}

%===============================================
\section{TEST DE INDEPENDENCIA}
%===============================================

\subsection{Tabla de contingencia}
Datos bivariados categóricos en tabla $I\times J$.

$H_0$: Las variables $A$ y $B$ son independientes.\\
$H_1$: Las variables están relacionadas.

\begin{formulabox}
$\chi^2=\sum_{i=1}^{I}\sum_{j=1}^{J}\frac{(O_{ij}-E_{ij})^2}{E_{ij}}\sim\chi^2_{(I-1)(J-1)}$\\[5pt]
$E_{ij}=\frac{n_{i\bullet}\cdot n_{\bullet j}}{n}$ (frec. esperada bajo independencia)
\end{formulabox}

\textbf{Condiciones}: $E_{ij}\geq 5$ para todas las celdas.

\R\texttt{chisq.test(tabla.contingencia)}

\subsection{Test de homogeneidad}
Misma fórmula. Contrasta si la distribución de una variable es igual en varias poblaciones.

%===============================================
\section{REGRESIÓN LINEAL SIMPLE}
%===============================================

\subsection{Modelo}
$Y_i=\beta_0+\beta_1 X_i+\varepsilon_i$ donde $\varepsilon_i\stackrel{iid}{\sim}N(0,\sigma_\varepsilon)$

\textbf{Hipótesis del modelo}:
\begin{itemize}[leftmargin=*]
\item Linealidad: relación lineal entre $X$ e $Y$
\item Independencia: errores independientes
\item Homocedasticidad: varianza constante de errores
\item Normalidad: errores normales
\end{itemize}

\subsection{Estimadores por mínimos cuadrados}
\begin{formulabox}
$b_1=\frac{\tilde{s}_{xy}}{\tilde{s}_x^2}=\frac{\sum_{i=1}^n(x_i-\bar{x})(y_i-\bar{y})}{\sum_{i=1}^n(x_i-\bar{x})^2}$\\[5pt]
$b_0=\bar{y}-b_1\bar{x}$
\end{formulabox}

\textbf{Fórmulas alternativas} (para cálculo directo con sumatorios):
\begin{formulabox}
$b_1=\frac{n\sum x_i y_i - \sum x_i \sum y_i}{n\sum x_i^2-(\sum x_i)^2}$\\[5pt]
$b_0=\frac{\sum y_i - b_1\sum x_i}{n}=\bar{y}-b_1\bar{x}$
\end{formulabox}

\textbf{Recta de regresión}: $\hat{y}=b_0+b_1 x$

\textbf{Residuos}: $e_i=y_i-\hat{y}_i$

\subsection{Propiedades}
\begin{itemize}[leftmargin=*]
\item La recta pasa por $(\bar{x},\bar{y})$
\item $\sum_{i=1}^n e_i=0$ (media de residuos = 0)
\item $\sum_{i=1}^n x_i e_i=0$
\item $b_0,b_1$ son estimadores insesgados de $\beta_0,\beta_1$
\end{itemize}

\subsection{Descomposición de variabilidad}
\begin{formulabox}
$SS_T=SS_R+SS_E$\\[3pt]
$\sum(y_i-\bar{y})^2=\sum(\hat{y}_i-\bar{y})^2+\sum(y_i-\hat{y}_i)^2$
\end{formulabox}
\begin{itemize}[leftmargin=*]
\item $SS_T$: suma de cuadrados total
\item $SS_R$: suma de cuadrados de regresión (explicada)
\item $SS_E$: suma de cuadrados del error (residual)
\end{itemize}

\subsection{Coeficiente de determinación}
\begin{formulabox}
$R^2=\frac{SS_R}{SS_T}=1-\frac{SS_E}{SS_T}=r_{xy}^2$
\end{formulabox}

\begin{itemize}[leftmargin=*]
\item $0\leq R^2\leq 1$
\item $R^2$ cercano a 1: buen ajuste
\item $R^2=$ proporción de variabilidad de $Y$ explicada por $X$
\end{itemize}

\subsection{Coeficiente de correlación lineal}
\begin{formulabox}
$r_{xy}=\frac{\tilde{s}_{xy}}{\tilde{s}_x\cdot\tilde{s}_y}=\frac{\sum(x_i-\bar{x})(y_i-\bar{y})}{\sqrt{\sum(x_i-\bar{x})^2}\sqrt{\sum(y_i-\bar{y})^2}}$
\end{formulabox}
$-1\leq r_{xy}\leq 1$. \quad $r_{xy}=\pm 1$: perfecta correlación lineal.

\subsection{Estimación de $\sigma_\varepsilon^2$}
\begin{formulabox}
$S^2=\frac{SS_E}{n-2}=\frac{\sum e_i^2}{n-2}$ (estimador insesgado)
\end{formulabox}

\subsection{Contraste sobre $\beta_1$}
$H_0:\beta_1=0$ (no hay relación lineal)
\begin{formulabox}
$T=\frac{b_1}{\text{SE}(b_1)}=\frac{b_1}{S/\sqrt{\sum(x_i-\bar{x})^2}}\sim t_{n-2}$
\end{formulabox}

p-valor $=2P(t_{n-2}>|t_0|)$

\textbf{Interpretación}:
\begin{itemize}[leftmargin=*]
\item $p<0.05$: regresión significativa, $\beta_1\neq 0$
\item $p>0.1$: no hay evidencia de relación lineal
\end{itemize}

\subsection{R: Regresión lineal}
\begin{verbatim}
modelo <- lm(Y ~ X, data=datos)
summary(modelo)
\end{verbatim}

\textbf{Elementos del summary}:
\begin{itemize}[leftmargin=*]
\item \texttt{Coefficients}: $b_0$, $b_1$, errores estándar, t-valor, p-valor
\item \texttt{Residual standard error}: $S$
\item \texttt{Multiple R-squared}: $R^2$
\item \texttt{F-statistic}: test global del modelo
\end{itemize}

\textbf{Otras funciones}:
\begin{itemize}[leftmargin=*]
\item \texttt{coef(modelo)}: coeficientes
\item \texttt{residuals(modelo)}: residuos
\item \texttt{fitted(modelo)}: valores ajustados $\hat{y}_i$
\item \texttt{predict(modelo,newdata=...)}: predicción
\item \texttt{confint(modelo)}: IC para coeficientes
\end{itemize}

%===============================================
\section{GUÍA RÁPIDA R}
%===============================================

\subsection{Distribuciones}
\textbf{Prefijos}: \texttt{d} densidad/prob., \texttt{p} distrib., \texttt{q} cuantil, \texttt{r} aleatorio

\begin{tabular}{@{}ll@{}}
\texttt{norm(x,$\mu$,$\sigma$)} & Normal\\
\texttt{t(x,df)} & t de Student\\
\texttt{chisq(x,df)} & Chi-cuadrado\\
\texttt{f(x,df1,df2)} & F de Fisher\\
\texttt{binom(k,n,p)} & Binomial\\
\texttt{pois(k,$\lambda$)} & Poisson\\
\texttt{geom(k,p)} & Geométrica\\
\texttt{nbinom(k,r,p)} & Binomial Negativa\\
\texttt{hyper(x,m,n,k)} & Hipergeométrica\\
\texttt{unif(x,a,b)} & Uniforme\\
\texttt{exp(x,rate)} & Exponencial
\end{tabular}

\subsection{Tests estadísticos}
\begin{tabular}{@{}ll@{}}
\texttt{t.test()} & Media(s)\\
\texttt{var.test()} & Dos varianzas (F-test)\\
\texttt{prop.test()} & Proporción(es)\\
\texttt{binom.test()} & Proporción exacto\\
\texttt{chisq.test()} & Bondad ajuste / Independencia\\
\texttt{fisher.test()} & Fisher exacto (2 proporciones)\\
\texttt{mcnemar.test()} & 2 proporciones emparejadas\\
\texttt{fligner.test()} & 2 varianzas (no paramétrico)\\
\texttt{sigma.test()} & Una varianza (TeachingDemos)\\
\texttt{varTest()} & Una varianza (EnvStats)\\
\end{tabular}

\subsection{Parámetros comunes en tests}
\begin{tabular}{@{}ll@{}}
\texttt{alternative} & \texttt{"two.sided"}, \texttt{"less"}, \texttt{"greater"}\\
\texttt{conf.level} & Nivel de confianza (0.95 por defecto)\\
\texttt{mu} & Media a contrastar (t.test)\\
\texttt{var.equal} & ¿Varianzas iguales? (t.test)\\
\texttt{paired} & ¿Muestras emparejadas? (t.test)\\
\texttt{p} & Proporción a contrastar\\
\end{tabular}

\subsection{Cuantiles más usados}
\begin{center}
\begin{tabular}{c|ccccc}
Confianza & 80\% & 90\% & 95\% & 98\% & 99\%\\
\hline
$\alpha$ & 0.20 & 0.10 & 0.05 & 0.02 & 0.01\\
$z_{1-\alpha/2}$ & 1.282 & 1.645 & 1.960 & 2.326 & 2.576
\end{tabular}
\end{center}

\subsection{Paquetes útiles}
\begin{itemize}[leftmargin=*]
\item \texttt{epitools}: intervalos para proporciones
\item \texttt{TeachingDemos}: \texttt{sigma.test}
\item \texttt{EnvStats}: \texttt{varTest}
\item \texttt{boot}: bootstrap
\end{itemize}

\subsection{Fórmulas de cálculo rápido}
\begin{formulabox}
\textbf{Media}: $\bar{x}=\frac{\sum x_i}{n}$\\[3pt]
\textbf{Varianza muestral}: $\tilde{s}^2=\frac{\sum x_i^2-n\bar{x}^2}{n-1}$\\[3pt]
\textbf{Covarianza}: $\tilde{s}_{xy}=\frac{\sum x_iy_i-n\bar{x}\bar{y}}{n-1}$\\[3pt]
\textbf{Correlación}: $r=\frac{\tilde{s}_{xy}}{\tilde{s}_x\tilde{s}_y}$\\[3pt]
\textbf{Pendiente regresión}: $b_1=\frac{\tilde{s}_{xy}}{\tilde{s}_x^2}=r\frac{\tilde{s}_y}{\tilde{s}_x}$
\end{formulabox}

\end{multicols}
\end{document}